% Gerado de docs/monografia/00-pre-textuais.md por scripts/monografia_latex.py.
% Não edite aqui: edite o .md e gere de novo.

Pointers are a difficult topic for novices learning C programming, and part of the difficulty stems not from a lack of knowledge but from misconceptions: ideas that are stable and coherent to those who hold them and that resist the correct explanation. A wrong answer alone does not distinguish a student who lacks the knowledge from one who holds a misconception, and the two cases call for different interventions. This work presents Structura, a system for teaching pointers in C that formulates, for each premise of the content, a diagnostic hypothesis from three signals provided by the student on a single screen (the answer, the reason behind it and the declared confidence) and leads each case along its own path: for the student who does not know, a sequence of short steps with immediate feedback; for the student who holds a misconception, a demonstration in which the student predicts the behavior of a program and then observes what the machine does (the program's output or the compiler's message); for the student who still errs after the steps, a review with the code animated line by line. The content covers 13 topics, with 79 premises, a catalog of 61 misconceptions with declared provenance, 402 steps, 80 demonstrations and 79 animated scenes. In the demonstrations, the scenes and 246 of the 402 steps, verification of the claims about program behavior does not rely on the judgment of those who wrote the material: an automatic checker compiles that code and, when the claim depends on execution, runs it, in the specified environment, and checks what the text claims, including the state represented in each scene; in the other 156 steps, which ask the student to choose among sentences, the code shown in 147 of them was checked only by reading; ten simulated student profiles, defined by fixed rules, go through every topic to check the specified behavior. The work does not evaluate the system with students: whether the diagnostic hypothesis matches the misconception a student actually holds and whether the content is pedagogically adequate depend on validation with people, whose protocol is presented as future work. It concludes that such a system can be built with verifiable behavior and with content that is verifiable as far as the compiler reaches, and that automated verification and the reading of the content identify defects of different kinds and are complementary.

\vspace{\onelineskip}
\noindent
\textbf{Keywords:} pointers in C; misconceptions; diagnosis; programming education; compiler-based verification.
